% Experimental setup tables for reproducibility (dt / corrector / Student-t VSC).
% Compile: pdflatex experimental_setup_reproducibility.tex
\documentclass[11pt]{article}
\usepackage[margin=1in]{geometry}
\usepackage{booktabs}
\usepackage{tabularx}
\usepackage{longtable}
\usepackage{multirow}
\usepackage{amsmath,amssymb}
\usepackage{hyperref}
\usepackage{caption}
\usepackage{enumitem}
\usepackage{array}

\newcommand{\trajid}{\texttt{traj\_id}}
\newcommand{\valmod}{\texttt{val\_mod}=5}

\title{Experimental Setup for Reproducibility\\
\large dt Head, $\varepsilon$-Corrector, and Student-$t$ VSC}
\author{}
\date{}

\begin{document}
\maketitle

\begin{abstract}
This note consolidates the missing reproducibility details requested by reviewers:
number of trajectories for dt-head and corrector training, train/validation splitting,
the number of residual samples used for Student-$t$ fitting, and whether a separate
corrector is trained per model / dataset / bitwidth.
Numbers follow the default scripts in this repository.
\end{abstract}

\section{High-level policy}

\begin{table}[t]
\centering
\caption{Answers to the main reproducibility questions.}
\label{tab:policy}
\begin{tabular}{@{}p{0.42\linewidth}p{0.52\linewidth}@{}}
\toprule
\textbf{Question} & \textbf{Answer} \\
\midrule
Separate corrector per model / dataset / bitwidth?
& \textbf{Yes.} Closed-loop trajectories are collected under the corresponding quantized sampler; correctors (and optional dt / mean heads) are \emph{not} shared across settings. \\
\addlinespace
Do dt and corrector share the same data?
& \textbf{No.} The dt head uses open-loop (or dt-labeled) trajectories; the $\varepsilon$-corrector / mean head uses closed-loop trajectories after quantization (optionally after dt); Student-$t$ fits use \emph{post-correction} residuals. \\
\addlinespace
Train / validation split?
& Trajectory-level: $(\trajid \bmod 5)=0$ $\rightarrow$ validation ($\approx 20\%$); otherwise training ($\approx 80\%$). We do \emph{not} shuffle individual timesteps across trajectories, to avoid leakage. \\
\bottomrule
\end{tabular}
\end{table}

\paragraph{Recommended manuscript sentence.}
\emph{For each tuple $(\mathrm{backbone},\mathrm{dataset},\mathrm{bitwidth})$, we collect closed-loop trajectories under that quantized sampler and train a dedicated corrector; checkpoints are not transferred across bitwidths or datasets.}

\section{Master configuration table}

\begin{table}[t]
\centering
\caption{Master protocol by backbone, dataset, and bitwidth.}
\label{tab:master}
\footnotesize
\begin{tabular}{@{}llllrrr@{}}
\toprule
\textbf{Backbone} & \textbf{Dataset} & \textbf{Quant.} & \textbf{Sampling} & \textbf{dt traj.} & \textbf{Corr.\ traj.} & \textbf{$t$-fit traj.} \\
\midrule
DDPM UNet & CIFAR-10 & W8A8
  & DDIM-100 (quad$\approx$96), $\eta{=}1$
  & (joint dt exp.) & mean: 500 & 200 \\
DDPM UNet & CIFAR-10 & W4A\{8,7,6,5\}
  & same
  & reuse W8A8 dt & mean: 200 / bit & 200 / bit \\
LDM UNet & ImageNet-256 & W4A8
  & DDIM-200, $\eta{=}1$, CFG$=$3.0
  & --- (no dt) & $\varepsilon$-corr.: 128 & 128 \\
LDM UNet & LSUN Bed./Church & W4A8 etc.
  & DDIM-200, $\eta{=}1$, no CFG
  & optional & typically 128 & same collect \\
U-ViT-S & CIFAR-10 & W8A8 / W4A8
  & DDIM-100 quad, $\eta{=}1$
  & open-loop 1000 & mean: 1000; VSC: 200 & 200 \\
\bottomrule
\end{tabular}
\end{table}

\section{dt-head training}

\begin{table}[t]
\centering
\caption{dt-head training settings.}
\label{tab:dt}
\begin{tabular}{@{}lll@{}}
\toprule
\textbf{Item} & \textbf{CIFAR UNet / U-ViT} & \textbf{ImageNet cin256 (corr.+tvar)} \\
\midrule
Used? & Yes & No (learned $\varepsilon$-corrector only) \\
Trajectory source & Open-loop quantized traj. & --- \\
\#trajectories & 1000 (U-ViT protocol) & --- \\
Steps / schedule & 100, \texttt{skip\_type=quad} ($\approx$96 updates) & --- \\
Labels & \texttt{label\_dt\_eps\_oracle} (e.g.\ \texttt{update\_mse}) & --- \\
Split & $(\trajid \bmod 5)=0$ $\rightarrow$ val (\valmod) & --- \\
Approx.\ train / val traj. & $\approx 800$ / $200$ & --- \\
Epochs / batch & 20 / 256 & --- \\
Checkpoint & \texttt{ckpt\_best.pt} by validation loss & --- \\
\bottomrule
\end{tabular}
\end{table}

\section{$\varepsilon$-corrector / mean residual training}

\begin{table}[t]
\centering
\caption{$\varepsilon$-corrector and mean residual-head training.}
\label{tab:corrector}
\footnotesize
\begin{tabular}{@{}lcccc@{}}
\toprule
\textbf{Item}
  & \textbf{ImageNet W4A8}
  & \textbf{CIFAR mean (W8A8)}
  & \textbf{CIFAR mean (W4A*)}
  & \textbf{U-ViT mean} \\
\midrule
Module
  & \texttt{DeltaEpsNet}
  & time-corr.\ residual net
  & same
  & same \\
\#trajectories
  & 128
  & 500
  & 200
  & 1000 (dt-only CL) \\
Steps / traj.
  & 200
  & $\approx$96
  & $\approx$96
  & $\approx$96 \\
\#states
  & $128\times200=25600$
  & $500\times96=48000$
  & $200\times96=19200$
  & $\approx 96000$ \\
Train / val traj.\ ($\approx$)
  & $102$ / $26$
  & $400$ / $100$
  & $160$ / $40$
  & $800$ / $200$ \\
Train / val states ($\approx$)
  & $20480$ / $5120$
  & $38400$ / $9600$
  & $15360$ / $3840$
  & proportional \\
Epochs
  & 20
  & 20
  & 20
  & 20 \\
Batch size
  & 8
  & 128
  & 128
  & script default \\
LR / optim.
  & $10^{-3}$, AdamW, cosine
  & same
  & same
  & same \\
Loss
  & MSE$+$cos$+$SR ($\lambda{=}1/2/0.3$)
  & MSE$+0.1\cdot$cos
  & same
  & same \\
Selection
  & best \texttt{val\_loss}
  & same
  & same
  & same \\
Seed
  & 1234
  & 1234
  & 1234
  & 1234 \\
\bottomrule
\end{tabular}
\end{table}

\paragraph{Per-bitwidth rule (CIFAR W4A*).}
Scripts such as \texttt{PTQD/vsc\_tvar/run\_w4a\_vsc\_tvar\_50k.sh} train a \emph{separate} mean checkpoint for each of W4A8 / W4A7 / W4A6.
ImageNet uses an independent directory \texttt{cin256\_w4a8\_.../learned\_corr/}.

\section{Student-$t$ variance fitting (VSC \texttt{var\_mle})}

\begin{table}[t]
\centering
\caption{Student-$t$ fitting protocol for VSC.}
\label{tab:studentt}
\begin{tabular}{@{}p{0.38\linewidth}p{0.56\linewidth}@{}}
\toprule
\textbf{Item} & \textbf{Setting} \\
\midrule
Input residual
  & Post-corrector residual $r=\varepsilon_{\mathrm{fp}}-\varepsilon_{\mathrm{corr}}$ \\
\#trajectories
  & Same as the traj.\ file used for fitting: ImageNet 128; CIFAR VSC estimate typically 200 \\
Aggregation
  & Independent fit at each nominal DDIM time index $t$ \\
Elements before subsample (CIFAR, 200 traj.)
  & $\approx N_{\mathrm{traj}}\times 3\times 32\times 32$ ($\approx 6.1\times 10^{5}$ per $t$) \\
Elements before subsample (ImageNet, 128 traj.)
  & $128\times 3\times 64\times 64$ ($\approx 1.6\times 10^{6}$ per $t$) \\
Subsample cap
  & \texttt{student\_t\_subsample}${}={}$\textbf{8000} elements per $t$ \\
Subsample seed
  & \texttt{student\_t\_seed}${}={}$0 \\
Minimum $n$
  & Skip MLE if $n<32$ \\
Estimator
  & \texttt{scipy.stats.t.fit} $\rightarrow (\nu,\sigma)$;
    $\mathrm{var\_mle}=\dfrac{\nu}{\nu-2}\sigma^{2}$ only if $\nu>2$ \\
Auxiliary statistic
  & 10\% two-sided trimmed mean of per-sample residual MSE (budget audit) \\
VSC inference
  & \texttt{--vsc\_var\_field var\_mle}, absorb${=}1.0$, \texttt{max\_budget\_fraction}${=}0.9$ \\
\bottomrule
\end{tabular}
\end{table}

\paragraph{Recommended manuscript sentence.}
\emph{Student-$t$ parameters are fit independently at each DDIM time index on up to 8{,}000 residual elements subsampled from the post-correction trajectories (seed 0); the MLE variance $\nu\sigma^{2}/(\nu-2)$ is stored in the VSC table.}

\section{Data flow}

\begin{enumerate}[leftmargin=*]
  \item Quantized model (per bitwidth).
  \item \emph{Optional:} open-loop trajectories $\rightarrow$ label $\mathrm{dt}^{\star}$ $\rightarrow$ train dt head.
  \item Closed-loop trajectories ($N_{\mathrm{traj}}\times T$).
  \item Train $\varepsilon$-corrector / mean head with split $(\trajid \bmod 5)=0$ for validation.
  \item Apply corrector $\rightarrow$ post-correction residuals.
  \item Per-time Student-$t$ MLE ($\le 8000$ elements) $\rightarrow$ VSC table $\rightarrow$ sampling / FID.
\end{enumerate}

\section{Hyperparameter checklist (appendix-ready)}

\begin{table}[t]
\centering
\caption{Default hyperparameters used across experiments.}
\label{tab:hparams}
\begin{tabular}{@{}ll@{}}
\toprule
\textbf{Hyperparameter} & \textbf{Default} \\
\midrule
Trajectory split & $\mathrm{val}=(\trajid \bmod 5)=0$ ($\approx 20\%$) \\
Corrector epochs & 20 \\
Corrector batch (CIFAR / ImageNet) & 128 / 8 \\
Corrector learning rate & $1\times 10^{-3}$ \\
dt epochs / batch & 20 / 256 \\
Student-$t$ per-step cap & 8000 elements \\
Student-$t$ subsample seed & 0 \\
FID sampling seed & 1234 \\
ImageNet CFG scale & 3.0 \\
VSC absorb / budget cap & $1.0$ / $0.9$ \\
\bottomrule
\end{tabular}
\end{table}

\section*{Implementation pointers}
\begin{itemize}[leftmargin=*]
  \item Corrector split: \texttt{noise\_eps\_corr/learned\_noise\_corrector.py} (\texttt{TrajectoryLateDataset}, \valmod).
  \item Corrector training: \texttt{noise\_eps\_corr/scripts/train\_corrector.py}.
  \item Mean residual (after dt): \texttt{state\_aware\_temporal\_joint/train\_dt\_corrected\_residual.py}.
  \item Student-$t$ VSC table: \texttt{PTQD/estimate\_time\_variance.py}.
  \item ImageNet pipeline: \texttt{PTQD/imagenet256/run\_cin256\_learned\_corr\_tvar.sh} ($N_{\mathrm{traj}}{=}128$, $T{=}200$).
\end{itemize}

\end{document}
